\section{Why the tagged divisor vector has no boundary}

The positive scalar law might suggest that the renormalized divisor series
converges in the original HCS-P51 Banach space.  The source tags rule this out.
Define
\begin{equation}
E_\tau=\tau^2\sum_{n\ge3}e^{-\tau n}D_n\in\Btag.
\label{eq:Etau}
\end{equation}
Because all coefficients are positive and different $n$-packets have
disjoint support,
\begin{equation}
\|E_\tau\|_{\rm tag}=\tau^2Z(\tau)\longrightarrow A_L>0.
\label{eq:Enorm}
\end{equation}

\begin{theorem}[Tagged-mass escape]
The family $(E_\tau)_{\tau>0}$ has neither a norm-convergent subnet nor a
weakly convergent subnet in $\Btag$ as $\tau\downarrow0$.
\end{theorem}

\begin{proof}
Fix a source coordinate $(\gamma_4,n,p)$.  Its coefficient in
\eqref{eq:Etau} is
\[
\tau^2e^{-\tau n}v_p(\beta_n)\longrightarrow0.
\]
Thus every coordinate of any norm limit would vanish.  The only possible
norm limit is zero, contradicting \eqref{eq:Enorm}.

The same coordinate functionals force any weak limit to be zero.  However,
the positive mass functional
\[
\mathfrak m\!\left(\sum c_{n,p}[\gamma_4,n,p]\right)
=\sum c_{n,p}\log p
\]
is bounded with norm one and satisfies
$\mathfrak m(E_\tau)=\|E_\tau\|_{\rm tag}\to A_L$.  Hence weak
convergence is also impossible.  Both contradictions apply to arbitrary
subnets.
\end{proof}

This theorem separates three data types:
\begin{center}
\begin{tabular}{lll}
\toprule
Object & Boundary behaviour & Status\\
\midrule
total packet mass & finite nonzero Abel constant & proved\\
scaled index distribution & $\Gamma(2,1)$ weak limit & proved\\
source-tagged divisor vector & no norm or weak subnet limit & refuted\\
\bottomrule
\end{tabular}
\end{center}
The first two rows cannot be promoted to the third.  Any boundary retaining
the full prime-ideal ledger must use a different topology or an additional
distributional construction.
